\documentclass[conference]{IEEEtran}
\IEEEoverridecommandlockouts

\usepackage{cite}
\usepackage{amsmath,amssymb,amsfonts}
\usepackage{graphicx}
\usepackage{textcomp}
\usepackage{xcolor}
\usepackage{booktabs}
\usepackage{multirow}
\usepackage{microtype}

\def\BibTeX{{\rm B\kern-.05em{\sc i\kern-.025em b}\kern-.08em
    T\kern-.1667em\lower.7ex\hbox{E}\kern-.125emX}}

\begin{document}

\title{MinImageScorer: Error-Driven Difficulty Scoring\\
for Guided Copy-Paste Augmentation in\\
Agricultural Object Detection}

\author{
\IEEEauthorblockN{[FILL: Author Name\textsuperscript{1},
                  Supervisor Name\textsuperscript{1}]}
\IEEEauthorblockA{%
[FILL: Department of XXX, University of XXX, Country]}
}

\maketitle

% ----------------------------------------------------------
\begin{abstract}
We propose \textbf{MinImageScorer} (MIS), an error-driven framework
that scores per-object and image difficulty by combining model
confidence and localization penalties on training-set failure signals.
Validated on MinneApple and Global Wheat Head Detection (GWHD), MIS
reliably ranks objects and images by difficulty and identifies
domain- and size-level failure patterns without additional information.
Score-guided copy-paste outperforms random selection on MinneApple
($\Delta\mathrm{AP}{=}{+}0.013$), but both strategies fall below
baseline. On GWHD, domain shift renders copy-paste ineffective
regardless of selection strategy, establishing MIS as a reliable
diagnostic for structural difficulty in agricultural detection.
\end{abstract}

\begin{IEEEkeywords}
object detection, copy-paste augmentation, difficulty scoring,
agricultural detection, small objects, domain shift
\end{IEEEkeywords}

% ----------------------------------------------------------
\section{Introduction}

Agricultural detection benchmarks present structural challenges absent
from general-purpose datasets. MinneApple~\cite{hani2020minneapple}
is dominated by densely packed sub-resolution apple instances where
failure is driven primarily by scale. GWHD~\cite{david2021global}
aggregates images from seven geographic domains, causing systematic
performance drops on minority-domain images. Copy-paste augmentation
yields gains on MS COCO~\cite{kisantal2019augmentation,ghiasi2021simple},
but applied blindly to structurally constrained datasets it risks
compounding failure modes: pasting sub-resolution objects amplifies
scale noise; pasting cross-domain objects amplifies distributional
mismatch.

We introduce \textbf{MinImageScorer} (MIS), a per-object difficulty
score derived from model failure signals, and show that MIS-guided
copy-paste outperforms random selection on MinneApple while both
strategies remain below baseline. On GWHD, copy-paste degrades
performance regardless of strategy, demonstrating that the dataset's
structural failure mode---not selection quality---determines whether
copy-paste can succeed.

% ----------------------------------------------------------
\section{Methodology}

\subsection{Scoring Formulation}

Let $G(I)$ be ground-truth objects in image $I$ and
$\mathcal{P}(g){=}\{p:\mathrm{IoU}(p,g)\ge\tau\}$ the predictions
matched to $g$, with $c_p{=}\mathrm{conf}(p)$ and
$u_p{=}\mathrm{IoU}(p,g)$. The \emph{Object Difficulty Score} is:
%
\begin{equation}
S_{\mathrm{obj}}(g)\!=\!
\begin{cases}
  \min_{p\in\mathcal{P}(g)}\!\bigl[\alpha(1{-}c_p){+}\beta(1{-}u_p)\bigr],
    & \mathcal{P}(g)\neq\varnothing,\\[2pt]
  \alpha{+}\beta, & \text{otherwise,}
\end{cases}
\label{eq:obj}
\end{equation}
%
where $\alpha,\beta$ weight confidence and localization penalties.
Missed objects receive maximum score $\alpha{+}\beta$.
The \emph{Image Difficulty Score} is:
%
\begin{equation}
S_{\mathrm{img}}(I)
  =w_1\!\cdot\!\frac{1}{|G(I)|}\!\sum_{g\in G(I)}\!S_{\mathrm{obj}}(g)
  +w_2\!\cdot\!\mathrm{FPrate}(I).
\label{eq:img}
\end{equation}
%
Weights are calibrated by maximizing Pearson correlation between
$S_{\mathrm{img}}$ and per-image miss and FP rates on the training
set, requiring no domain-specific prior knowledge.

\subsection{Score-Guided Copy-Paste Pipeline}

A baseline YOLOv8s detector is trained, then full-training-set
inference yields per-prediction confidence and IoU values from which
MIS scores are computed. The dataset is extended by 30\%, with 2--8
objects pasted per new image using scale jitter and random
flipping~\cite{ghiasi2021simple}. Under \emph{Score CP}, source
objects are sampled proportionally to $S_{\mathrm{obj}}$ via
exponential weighting; under \emph{Random CP}, uniformly.
Both conditions use identical hyperparameters.

% ----------------------------------------------------------
\section{Results and Discussion}

\subsection{Score Validation}

Table~\ref{tab:validation}(a) shows strong positive correlations
($r{>}0.69$) between $S_{\mathrm{img}}$ and per-image miss and FP
rates, confirming the score reliably captures image-level difficulty.
Table~\ref{tab:validation}(b) shows that harder GWHD domains (lower
AP) carry higher mean $S_{\mathrm{img}}$, matching the known
cross-domain failure mode. Table~\ref{tab:validation}(c) shows
$S_{\mathrm{obj}}$ is higher for small objects than medium objects on
MinneApple, correctly mirroring AP ($0.300$ vs.\ $0.592$). MIS
therefore recovers both domain- and size-level difficulty using only
training-set model failure signals.

\begin{table}[t]
\centering
\caption{Score validation. \textbf{(a)}~Pearson $r$ of $S_{\mathrm{img}}$
with miss/FP rates. \textbf{(b)}~GWHD domain AP vs.\ mean
$S_{\mathrm{img}}$. \textbf{(c)}~MinneApple scores by object size.}
\label{tab:validation}
\resizebox{\columnwidth}{!}{%
\begin{tabular}{@{}lcc@{}}
\toprule
\multicolumn{3}{@{}l}{\textbf{(a) Correlation with detection failures}}\\
\textbf{Dataset} & $r$(\,score,\,miss\,) & $r$(\,score,\,FP\,)\\
\midrule
MinneApple & 0.698 & 0.688 \\
GWHD       & 0.816 & 0.818 \\
\midrule
\multicolumn{3}{@{}l}{\textbf{(b) GWHD domain-level difficulty}}\\
\textbf{Domain} & \textbf{AP} & \textbf{Mean }$S_{\mathrm{img}}$\\
\midrule
rres\_1    & 0.572 & 0.200 \\
ethz\_1    & 0.550 & 0.227 \\
inrae\_1   & 0.523 & 0.251 \\
arvalis\_3 & 0.466 & 0.294 \\
usask\_1   & 0.408 & 0.312 \\
arvalis\_1 & 0.403 & 0.331 \\
arvalis\_2 & 0.365 & 0.424 \\
\midrule
\multicolumn{3}{@{}l}{\textbf{(c) MinneApple object scores by size}}\\
\textbf{Metric}
  & \textbf{Small ($\le$32\,px)}
  & \textbf{Medium ($\le$96\,px)}\\
\midrule
AP                       & 0.300 & 0.592 \\
$S_{\mathrm{obj}}$ (mean)& 0.378 & 0.222 \\
\bottomrule
\end{tabular}%
}
\end{table}

\subsection{Augmentation Results}

Table~\ref{tab:ap} summarises copy-paste outcomes. On MinneApple,
Score CP outperforms Random CP across all sub-metrics
($\Delta\mathrm{AP}{=}{+}0.013$), confirming that difficulty-weighted
sampling improves selection quality. Both strategies nonetheless fall
below baseline: boundary artifacts from unblended pastes reduce
localization precision, and sub-resolution objects cannot be made more
learnable through repetition alone. On GWHD, Score CP and Random CP
are statistically indistinguishable and both degrade below baseline.
Score-guided sampling concentrates objects from domain-minority
sources, amplifying distributional imbalance rather than correcting
it. These results confirm that covariate shift between augmented and
test images---not selection strategy---governs GWHD performance.

\begin{table}[t]
\centering
\caption{Copy-paste augmentation results.}
\label{tab:ap}
\resizebox{\columnwidth}{!}{%
\begin{tabular}{@{}llrrrrr@{}}
\toprule
\textbf{Dataset} & \textbf{Condition}
  & \textbf{AP} & \textbf{AP$_{50}$} & \textbf{AP$_{75}$}
  & \textbf{AP$_S$} & \textbf{AP$_M$}\\
\midrule
\multirow{3}{*}{MinneApple}
  & Baseline          & 0.439 & 0.771 & 0.452 & 0.300 & 0.592 \\
  & Random CP         & 0.415 & 0.741 & 0.427 & 0.279 & 0.563 \\
  & \textbf{Score CP} & \textbf{0.428} & \textbf{0.759}
                      & \textbf{0.440} & \textbf{0.286} & \textbf{0.581}\\
\midrule
\multirow{3}{*}{GWHD}
  & Baseline          & 0.507 & 0.922 & 0.491 & 0.136 & 0.501 \\
  & Random CP         & 0.486 & 0.905 & 0.456 & 0.118 & 0.480 \\
  & \textbf{Score CP} & \textbf{0.484} & \textbf{0.905}
                      & \textbf{0.455} & \textbf{0.098} & \textbf{0.477}\\
\bottomrule
\end{tabular}%
}
\end{table}

% ----------------------------------------------------------
\section{Conclusion}

MIS reliably quantifies per-object detection difficulty from model
failure signals, recovering size- and domain-level structural
patterns without privileged information. Score-guided copy-paste
outperforms random selection on MinneApple but degrades below
baseline on both datasets: sub-resolution scale limits and
cross-domain distributional shift are structural constraints that
augmentation alone cannot overcome. Copy-paste is beneficial when
objects are underrepresented yet visually compatible with the test
distribution; it is counterproductive under resolution limits or
domain shift. Future work should pursue better boundary blending to
reduce distributional displacement and dynamic score recomputation
to adapt augmentation as the model improves.

% ----------------------------------------------------------
\begin{thebibliography}{00}

\bibitem{hani2020minneapple}
N.~Hani, P.~Roy, and B.~Bhanu,
``MinneApple: A Benchmark Dataset for Apple Detection and Segmentation,''
\textit{IEEE Robotics and Automation Letters},
vol.~5, no.~2, pp.~852--858, 2020.

\bibitem{david2021global}
E.~David \textit{et~al.},
``Global Wheat Head Detection (GWHD) Dataset,''
\textit{Plant Phenomics}, vol.~2021, 2021.

\bibitem{kisantal2019augmentation}
M.~Kisantal \textit{et~al.},
``Augmentation for Small Object Detection,''
in \textit{Proc.\ ICACCI}, 2019.

\bibitem{ghiasi2021simple}
G.~Ghiasi \textit{et~al.},
``Simple Copy-Paste is a Strong Data Augmentation Method
for Instance Segmentation,''
in \textit{Proc.\ IEEE/CVF CVPR}, pp.~2918--2928, 2021.

\end{thebibliography}

\end{document}